\chapter{综合流行病学调查、微分方程理论和神经网络模型的方法} \label{ch:methods}


\section{数据驱动下的方法综合}

本节讨论了在数据驱动下的综合方法，以解决流行病学问题。主要涵盖了以下三个关键方面：

流行病学调查方法强调了数据的重要性，包括收集各种数据，如病例报告、感染人群的特征和传播时间线等。
这些数据用于描述流行病学特征、验证模型预测以及追踪疫情的变化。

微分方程理论在理解流行病的机制中起到关键作用。
微分方程被用来建立流行病模型，描述疫情的动态行为。
微分方程理论用于确定模型的平衡点性质，包括无病平衡点和地方病平衡点，以及阈值，如基本再生数等。

神经网络可以被引入以处理流行病学问题。
神经网络具有非线性建模能力、数据驱动性以及在预测性能方面的优势。
神经网络能够处理非线性效应，从数据中学习潜在模式，并用于预测流行病发展趋势。


\section{流行病学调查}

流行病学调查是流行病学研究的关键组成部分。
数据收集是流行病学调查的基础，各种数据，如病例报告、感染人群的特征以及疫情传播的时间线，都必不可少。
通过描述流行病学特征，以了解感染的地理分布、感染人群的年龄、性别和健康状况等信息，这些信息对于调整模型结构和进行参数估计非常重要。
验证模型预测是确保流行病模型准确性和适用性的关键步骤，通过比较模型的预测与实际疫情发展趋势，评估模型的性能。
流行病学调查还包括追踪疫情的变化，这包括监测新变种的出现和传播，以及疫情的季节性变化等，这些信息对于制定干预策略和应对流行病的挑战至关重要。


\section{微分方程理论}

微分方程被用来建立流行病学模型，描述疫情的动态行为。这些微分方程基于已知的流行病学参数和初始条件，描述了感染病例、康复病例、潜伏期病例等之间的关系。通过建立模型，研究人员可以更好地理解流行病学过程，预测疫情的发展趋势。

微分方程模型允许研究人员估计模型中的参数，例如感染率、康复率、潜伏期等。这些参数估计通常通过将模型的预测与实际观测数据进行拟合来获得。参数估计的准确性对于模型的可信度和适用性至关重要，因为它们反映了疫情的真实情况。

微分方程理论还用于分析模型的稳定性。这包括确定模型解是否会趋于稳定状态，这对于预测疫情的最终结果非常重要。通过分析平衡点的性质，包括存在唯一性、局部稳定性、全局稳定性和绝灭性等方面，可以评估模型的稳定性。


\section{神经网络模型}

神经网络模型在流行病学中具有重要作用，它能够进行非线性建模，有效捕捉复杂的流行病学动态，特别适用于模拟实际疫情中的非线性效应。
神经网络也是数据驱动的工具，可以从大规模流行病学数据集中学习潜在的模式和关联性，为深入了解流行病提供强有力的支持。
神经网络在时间序列预测和趋势分析方面表现出色，可用于准确预测未来的流行病发展趋势。
综合这些优势，神经网络模型为深入理解和有效应对流行病学问题提供了一种强大的工具。

受第~\ref{ch:introduction}中已有模型的启发，本文构建的 LSTM 模型结构如下：

\begin{algorithm}
	\label{alg:algorithm}
	Model(

	\qquad (lstm): LSTM(8, 32, batch\_first=True)

	\qquad (fc\_1): Linear(in\_features=32, out\_features=128, bias=True)

	\qquad (fc\_2): Sigmoid()

	\qquad (fc\_3): Tanh()

	\qquad (fc): Linear(in\_features=128, out\_features=2, bias=True)

	\qquad (relu): ReLU()

	)
\end{algorithm}

为了评估预测值与实际值之间的差异，找出梯度下降方向以减小间隙，根据以下方程，将该模型的损失函数设置为均方误差(MSE) :

\begin{equation*}
	\text{MSE}=\frac{\sum_{i=1}^n\left(y_i-\hat{y}_i\right)^2}{n}
\end{equation*}

设置 LSTM 模型使用Adam优化器、学习率设置为0.001、损失函数采用均方误差函数（MSE）、使用随机梯度下降法进行训练，并进行 200 次迭代。